
\section{Conclusiones}
\label{sec:conclusiones}

% -----------------------------------------------------------------------
\subsection{Síntesis del trabajo realizado}
\label{sec:conclusiones-sintesis}

MDP-ProbLog se extendió para representar fluentes de estado multivaluados
mediante Disyunciones Anotadas (ADs), un constructo nativo de la
Programación Lógica Probabilística que garantiza la exclusión mutua entre
los posibles valores de una variable por construcción semántica
\cite{vennekens2004lpads}. Bajo esta representación, el espacio de estados
adopta una estructura en base mixta donde cada fluente opera con su
cardinalidad real, lo que elimina los estados espurios generados por la
binarización forzada de la versión original y traslada la responsabilidad
de la exclusión mutua del modelador al motor lógico.

La intervención se limita a la capa de representación que alimenta al
compilador de ProbLog, sin alterar los algoritmos de inferencia ni los
circuitos compilados. Ello preserva la compatibilidad con todos los modelos
booleanos existentes y las garantías formales de Bueno et al.\
\cite{bueno2016mdp}.

La evaluación empírica sobre los dominios de Mitchell \cite{mitchell1997machine}
y Russell y Norvig \cite{russell2016artificial} confirmó que ambas
codificaciones producen la misma $V^*$ y la misma $\pi^*$ para un mismo
MDP, en condiciones deterministas y estocásticas. Los experimentos mostraron
además que la codificación multivaluada mantiene el tamaño del programa
fuente constante respecto al tamaño del dominio y elimina los estados
espurios. La mejora computacional resultante crece con la escala del
problema.

% -----------------------------------------------------------------------
\subsection{Contribuciones}
\label{sec:conclusiones-contribuciones}

Se aporta un esquema factorizado del espacio de estados bajo una
estructura \textit{mixed-radix} que asigna a cada fluente su cardinalidad
real. A diferencia de la representación booleana original, donde el tamaño
del espacio de estados crece como potencia de dos del número de bits
independientemente de la semántica del dominio, la representación
\textit{mixed-radix} garantiza que cada índice corresponda exactamente a un
estado válido. Esta propiedad elimina por construcción los estados espurios
que la codificación booleana introduce cuando el número de estados del
dominio no es una potencia exacta de dos \cite{hoey1999spudd}. Junto con el
esquema factorizado, se desarrolló un clasificador automático de fluentes
que determina, a partir de un índice invertido sobre los nodos
\texttt{choice} de la base de cláusulas, si un fluente declarado por el
usuario es booleano o multivaluado. El clasificador soporta tanto la
inferencia implícita del tipo, mediante el análisis de las ADs presentes en
el programa, como la declaración explícita mediante el predicado
\texttt{state\_fluent/2}.

La integración opera en la fase de inyección de hechos, donde
sustituye los hechos probabilísticos independientes por ADs uniformes para
los fluentes multivaluados. Ello preserva las garantías de inferencia exacta
sobre circuitos d-DNNF y SDD establecidas por Darwiche y Marquis
\cite{darwiche2002knowledge, fierens2015inference}. El sistema conserva
además la retrocompatibilidad con los modelos booleanos existentes. Cuando
el programa del usuario no contiene ADs, el \textit{pipeline} se comporta
de forma idéntica al de Bueno et al.\ \cite{bueno2016mdp}.

Se implementó también el evaluador de diferenciación de circuitos de
Darwiche \cite{darwiche2003differential}, que calcula todas las marginales
requeridas durante la Iteración de Valor en dos pasadas lineales sobre el
circuito compilado. Este evaluador reduce el costo de evaluar $Q$ consultas
simultáneas de $\mathcal{O}(Q \times |C|)$ a $\mathcal{O}(|C|) +
\mathcal{O}(Q)$, donde $|C|$ denota el tamaño del circuito. El beneficio es
multiplicativo con el de la codificación factorizada: aquella reduce el
número de pares estado-acción por iteración, mientras que este reduce el
costo de procesar cada par. El sistema admite además modelos mixtos en los
que fluentes booleanos y multivaluados coexisten en el mismo programa, lo
que permite una adopción gradual sin reescribir los modelos existentes.

% -----------------------------------------------------------------------
\subsection{Limitaciones técnicas identificadas}
\label{sec:conclusiones-limitaciones}

Durante el desarrollo e implementación se identificaron limitaciones técnicas
distintas de las restricciones de diseño anticipadas en la
Sección~\ref{subsec:limitaciones}.

\begin{enumerate}
    \item \textbf{Agrupamiento global de fluentes.} El clasificador agrupa
    bajo un único fluente multivaluado todos los términos que comparten
    predicado y aridad. Esta convención impide representar dominios en los
    que múltiples entidades independientes comparten el mismo predicado con
    valores distintos.

    \item \textbf{Incompatibilidad entre la configuración SDD y dominios
    estocásticos.} La configuración de inferencia \textit{sdd} generó un
    error de compilación en el dominio Russell, que presenta transiciones
    estocásticas no triviales. La causa no fue diagnosticada; se presume una
    restricción del \textit{backend} SDD de ProbLog2 ante ciertos patrones
    de ADs.

    \item \textbf{Supuesto de independencia condicional no verificado en
    tiempo de ejecución.} El cálculo del valor esperado factorizado asume
    que las distribuciones de transición de cada fluente son condicionalmente
    independientes dado el estado y la acción actuales. El sistema no
    detecta ni advierte si las reglas del usuario introducen dependencias
    entre fluentes que violan este supuesto, lo que puede producir
    resultados incorrectos sin aviso.
\end{enumerate}

\clearpage

% -----------------------------------------------------------------------
\subsection{Trabajo futuro}
\label{sec:conclusiones-futuro}

Las limitaciones identificadas y la estructura de la extensión propuesta
sugieren las siguientes líneas de desarrollo e investigación.

\begin{enumerate}
    \item \textbf{Fluentes con múltiples entidades.} La limitación de
    agrupamiento global impide modelar con naturalidad dominios en los que
    varias entidades independientes comparten el mismo predicado de fluente.
    Extender el clasificador para operar por entidad, diferenciando
    instancias a partir de sus argumentos de identificación, cubriría una
    clase más amplia de dominios sin requerir declaraciones adicionales del
    usuario.

    \item \textbf{Resolución simbólica del espacio de estados.} Esta
    extensión no modificó la estrategia de resolución enumerativa de
    MDP-ProbLog; el algoritmo de Iteración de Valor continúa recorriendo
    todos los pares estado-acción en cada iteración. Vale la pena explorar la
    viabilidad de adaptar ideas historicamente aplicadas a otro tipo de representaciones
    simbólicas, como los Diagramas de Decisión Algebraicos de SPUDD \cite{hoey1999spudd}
    y extraporalas a los Diagramas de Decisión Sentenciales o los d-DNNFs. 
    Esto permitiría manipular la función de valor sin
    enumerar el espacio de estados y superaría los límites de escalamiento
    observados en los experimentos.

    \item \textbf{Impacto de las Disyunciones Anotadas en el tamaño del
    circuito compilado.} Los experimentos mostraron que la codificación
    factorizada produce tiempos de resolución menores, pero el mecanismo por
    el cual las ADs afectan el tamaño del circuito d-DNNF o SDD no fue
    modelado teóricamente. Cuantificar esta relación permitiría predecir el
    comportamiento del compilador ante distintas estructuras de fluentes
    multivaluados y orientar decisiones de diseño de modelos.
\end{enumerate}
